\documentclass{article}
\usepackage[utf8]{inputenc}

\usepackage{amsthm}
\usepackage{amsmath}
\usepackage{amssymb}
\usepackage{enumitem}
\usepackage{ marvosym }
\usepackage{amscd}

\title{Product of Orders}
\author{Adam W. Grant}
\date{March 2021}

\newtheorem*{defn}{Defnition}

\begin{document}

\maketitle

\section{Introduction}
This short note presents a derivation of a pointfree version of the definition of the product of two partial orders.

\section{Product of Orders}
The following definition is from \cite{10.5555/551462}
\begin{defn}
Let $(A,\sqsupseteq_A)$ and $(B,\sqsupseteq_B)$ be two posets.  We define a partial ordering on the cartesian product $A \times B$ componentwise by
\begin{equation}
    (a',b') \sqsupseteq_{A\times B} (a,b) \;\equiv\; a' \sqsupseteq_A a \wedge b' \sqsupseteq_B b \tag{$\sqsupseteq$ of pairs}
\end{equation}
\end{defn}
It is this definition that we will subject to the Pointfree Transform\footnote{See \cite{Oliveira2009} for an explanation of the process.}, in what remains of this note.

\section{Removing the Points}
Here we calculate the pointfree equivalent of equation ($\sqsupseteq$ of pairs):
\renewcommand{\labelitemii}{\textbullet}
\begin{itemize}
    \item $(a',b') \sqsupseteq_{A\times B} (a,b)$
    \item[$\equiv$] \{ ($\sqsupseteq$ of pairs) \}
    \item[] $a' \sqsupseteq_A a \wedge b' \sqsupseteq_B b$
    \item[$\equiv$] \{ Definitions of $fst$ and $snd$ \}
    \item[] $fst(a',b') \sqsupseteq_A fst(a,b) \wedge snd(a',b') \sqsupseteq_B snd(a,b)$
    \item[$\equiv$] \{ Guardanapo, twice \}
    \item[] $(a',b') (fst^{\circ} \cdot \sqsupseteq_A \cdot fst) (a,b) \wedge (a',b') (snd^{\circ} \cdot \sqsupseteq_B \cdot snd) (a,b)$
    \item[$\equiv$] \{ Meet \}
    \item[] $(a',b') (fst^{\circ} \cdot \sqsupseteq_A \cdot fst \cap snd^{\circ} \cdot \sqsupseteq_B \cdot snd) (a,b)$
    \item[\Heart]
\end{itemize}
That is, we have:
\begin{equation*}
    \sqsupseteq_{A\times B} \;:=\; fst^{\circ} \cdot \sqsupseteq_A \cdot fst \cap snd^{\circ} \cdot \sqsupseteq_B \cdot snd
\end{equation*}
which is free of points, as desired.

\bibliographystyle{plain}
\bibliography{refs}

\end{document}
